\section{Implementierung eines Lichtgatters in das Gesamtkonzept}
%%%%%%%%%%%%%%%%%%%%%%%%%%%%%%%%%%%%%%%%%%%%%%%%%%%%%%%%%%%%%%%%%%%%%%%%%%%%%%%%%%%%%%%%%%%%%%%%%%%%%%%%%%%%%%%%%%%%%%%%

Im aktuellen Ist-Zustand des Hochregallagers fehlt die vollständige Implementierung
der Sicherheitslogik.
Lediglich der Not-Aus-Schalter ist hardwareseitig verdrahtet und schaltet bei Betätigung
die Zuleitung ab. Es existiert jedoch keine automatische Abschaltung im Gefahrenfall.
Die Logik des Ein- und Auslagerns ist somit nicht gefahrenfrei. Sobald ein Eingriff in den
Prozess erfolgt, kann dieser im aktuellen Zustand nicht automatisch gestoppt werden.
Da das Miniaturmodell eine reale Situation in der Logistik abbilden soll, stellt dieser
Punkt ein erhebliches Defizit dar.
Es gibt verschiedene Möglichkeiten, um diesen Nachteil zu beheben. Die naheliegendste
Option ist die Implementierung einer Lichtgatter-Funktion, um den Gefahrenbereich
zuverlässig zu überwachen und den Betrieb im Störfall sicher zu unterbrechen. Durch
ein vorheriges Projekt wurde bereits, ein Modell der Firma IFM[8], installiert.

\subsection{Zielsetzung der Implementierung}
Die Studienarbeiten der vergangenen Jahre haben das Hochregallager in einen Zustand
gebracht, der wichtige Funktionen wie Ein- und Auslagern abbildet.
Jedoch wurden hierbei, fälschlicherweise, die notwendigen Sicherheitsaspekte völlig
außer Acht gelassen.
In einer realen Umgebung, beispielsweise in Materiallagern mit Hochregalsystemen
großer Industrieunternehmen, wäre eine solche Vorgehensweise undenkbar.
Es muss jederzeit gewährleistet sein, dass Gefahren so gut wie möglich unter Kontrolle
gehalten werden. Das bedeutet, im Fehlerfall muss jede gefährliche Bewegung sicher
abgeschaltet werden.
Dies stellt die zentrale Zielsetzung der Implementierung dar. In jedem Betriebsmodus
muss sichergestellt sein, dass bei Erkennung einer Gefahr automatisch abgeschaltet
wird.
Dabei ist jedoch zu berücksichtigen, dass der Kolben des Zylinders, welcher die Mate
rialien einlagert, kein Abschaltsignal auslösen darf. Dieser fährt nämlich unmittelbar in
den Gatterbereich der Lichtschranke ein.
Somit ist zunächst eine umfassende Analyse des Ist-Zustands sowie der verschiedenen
Prozesswerte des Lichtgatters erforderlich.
Erst im Anschluss kann die Implementierung der Sicherheitslogik erfolgen

\subsection{Analyse der Paramerierung des Lichtgatters}
Die Konfiguration des Lichtgatters ist entscheidend für die richtige Funktionsweise. Beim
Modell von IFM[8] gibt es mehrere Möglichkeiten zur Parametrierung.
Das Lichtgatter besitzt mehrere Lichtstrahlen, welche über einzelne Bits definiert sind.
Sobald sich ein Objekt im Gatter befindet, werden mehrere dieser Lichtstrahlen gleich
zeitig unterbrochen.
Der Sensor des vorliegenden Lichtgatters hat im sogenannten Detektionsmodus
(GTBO) den Wert 1. Das bedeutet, dass er schaltet, sobald ein Objekt die Strahlen
unterbricht.
Über IO-Link kann nun festgelegt werden, ab wie vielen unterbrochenen Strahlen der
Sensor ein Schaltsignal ausgeben soll

\subsubsection{Prozesswerte der Steuerung des Arbeitsverhaltens}
Zur Steuerung des Verhaltens bei einer Unterbrechung stehen verschiedene Prozess
werte zur Verfügung[7]:
• FBO(First Beam Occupied)
• LBO(Last Beam Occupied)
• CBO(Central Beam Occupied)
• NBO(Number of Beams Occupied)
• NCBO(Number of Consecutive Beams Occupied)
Im Betriebsmodus FBO schaltet der Sensor bereits beim ersten unterbrochenen Licht
strahl. Das bedeutet, Objekte, die im unteren Bereich des Lichtgatters liegen, werden
als erstes erfasst.
Im Modus LBO geschieht genau das Gegenteil: Es wird beim letzten Lichtstrahl ge
schaltet.
Dieser Modus wird beispielsweise bei Endlagenerkennungen oder beim Auslagern
verwendet– also immer dann, wenn das Austreten aus einem bestimmten Bereich
überwacht werden soll.
Der Modus CBO ist etwas komplexer: Hier wird der mittlere Lichtstrahl (bei mehreren
zusammenhängenden Strahlen) überwacht. Tritt ein Objekt in das Lichtgatter ein, wird
der zentrale Strahl als Referenzpunkt betrachtet. Bei mehreren Objekten wird immer
vom größeren Objekt ausgegangen, also von dem, das mehr Strahlen gleichzeitig
unterbricht– wie in Abbildung 3.13 zu sehen.
Automatisiertes Modell- Hochregallager
26
3 Implementierung eines Lichtgatters in das Gesamtkonzept
Anwendungsbeispiele sind z. B. Durchgänge oder Förderabschnitte, also überall dort,
wo Objekte mittig erfasst werden sollen (Referenzpunkt: Mitte).
Abbildung 3.13: Modus CBO; Quelle: [7]
Die letzten beiden Modi, NBO und NCBO, beschäftigen sich mit der konkreten Anzahl
der Lichtstrahlen, die unterbrochen werden.
Der Unterschied der beiden Modi liegt darin, dass NBO die Gesamtanzahl der unterbro
chenen Lichtstrahlen angibt, während NCBO die Anzahl der Strahlen beschreibt, die
aufeinanderfolgend unterbrochen werden.
Im Projekt wird der Modus NBO verwendet. Eine Abbildung soll die Funktionsweise
noch einmal verdeutlichen (siehe Abbildung 3.14)

\subsection{Komponentenübersicht und Systemaufbau}
FürdieUmsetzungderSicherheitsfunktionwurdeeinSystemaufbaurealisiert,derdie
LichtschrankedirektmiteinemIO-Link-MasterundderSPSverbindet.ImFolgenden
werdendieeingesetztenKomponentensowiederenZusammenspiel innerhalbdes
Gesamtsystemsbeschrieben.
Komponente Hersteller/Typ Funktion Schnittstelle
Lichtgitter / Licht
schranke
IFMOY5100(Infrarot) Erfassung von Objekten
undÜberwachungdesGe
fahrenbereichs
IO-Link
IO-LinkMaster IFMAL1100 Kommunikation zwischen
IO-Link-Devices und der
SPS
PROFINET /
IO-Link
SPS SiemensS7-1500 Steuerungdesgesamten
Systems,Verarbeitungder
Sensorsignale
PROFINET
Netzteil24VDC SiemensSITOP/ver
gleichbar
Versorgung von IO-Link
Master,Lichtschrankeund
SPS
24VDC
Aktorik(Zylinder/Mo
tor)
Festo(pneumatisch) Ausführungder Ein- und
Auslagerbewegung
Digitale Aus
gänge der
SPS
Tabelle4:ÜbersichtderverwendetenKomponentenimSystemaufbau
InTabelle4sinddieeinzelnenKomponentendesHochregallagersrundumdieLicht
schrankedargestellt.
Diesemüssenoptimalmiteinanderinteragieren,umdiegewünschteFunktionalitätsi
cherzustellen.
DieLichtschrankeistvomTypOY5100undstammtvomHerstellerIFM.Eshandeltsich
hierbeiumeinsteuerbaresLichtgitter,dasüberIO-LinkmitderSPSkommuniziert.

\subsubsection{Systemaufbau}
Die folgende Abbildung zeigt den strukturellen Aufbau des Gesamtsystems in verein
fachter Form.
SPS
(Siemens S7-1500)
Netzteil
24V DC
24V DC
PROFINET
IO-Link Master
(IFM AL1100)
IO-Link
Lichtgitter
(IFM OY5100)
Signal / Reaktion
Aktorik
(Zylinder / Motor)
Abbildung 3.15: Vertikaler Systemaufbau des Lichtgatter-Systems
Im Diagramm 3.15 ist der Aufbau der einzelnen Komponenten grafisch dargestellt. Der
IO-Link-Master dient hierbei als Schnittstelle zwischen SPS und Lichtgitter und steuert
die Kommunikation über das im TIA Portal implementierte Programm.
Im Fehlerfall werden gefährliche Bewegungen, die beispielsweise vom Zylinder oder
den Motoren der Linearachsen ausgehen, automatisch gestoppt.
Im folgenden Abschnitt wird die Implementierung dieser Sicherheitslogik in der SPS
detailliert erläutert

\subsection{Implementieren der Sicherheitslogik in die SPS}
Die Sicherheitslogik wird, wie bereits erwähnt, in eine Siemens S7-1500 SPS einpro
grammiert. Die Programmierung erfolgt mithilfe der Software TIA Portal.
Für die Umsetzung der Arbeit wurde das bestehende Programm des Jahrgangs TEA
22 als Grundlage verwendet. Zunächst wurde eine neue Program Organization Unit
(POU) mit der Bezeichnung Lichtschranke erstellt.
Im Anschluss erfolgte das Einbinden des Moduls OY5100 im Reiter Geräte und Netz
werke. Der Adressbereich (frei wählbar) wurde so konfiguriert, dass eine eindeutige
Zuordnung innerhalb des Projekts gewährleistet ist.
Abbildung 3.16: Einbindung der IFM-Lichtschranke OY5100 im TIA Portal
Die zugewiesenen Adressen werden für die verschiedenen Prozesswerte genutzt (siehe
Kapitel 3.2.1).
Als Vorbereitung auf die Programmierung wurde eine sogenannte Watchtable erstellt.
Diese dient dazu, die einzelnen Prozesswerte des Lichtgatters während der Laufzeit zu
überwachen und deren Verhalten zu analysieren.
Die Prozesswerte liegen im Format WORD (2 Byte) vor und bilden die Anzahl der
unterbrochenen Lichtstrahlen ab.
Nach einer Reihe von Tests wurde festgestellt, dass der Kolben des Zylinders im Ar
beitsbereich des Lichtgatters drei Strahlen gleichzeitig unterbricht. Hierfür erfolgte
ein gezieltes Verfahren der Lineareinheit über alle vier Ebenen des Hochregallagers.
So konnte überprüft werden, dass der Kolben in jeder Position konstant drei Strahlen
abdeckt.
Im weiteren Verlauf wurde diskutiert, welcher Prozesswert sich für die Sicherheitslogik
am besten eignet. Es zeigte sich, dass der Modus NBO (Number of Beams Occupied)
Automatisiertes Modell- Hochregallager
30
3 Implementierung eines Lichtgatters in das Gesamtkonzept
die Summe aller unterbrochenen Lichtstrahlen angibt (siehe Kapitel 3.2.1). Daher ist
dieser Prozesswert optimal geeignet, um eine zuverlässige Abschaltung über das Licht
gatter zu realisieren.
Nach der Diskussion wurde, für die Realisierung der Abfrage, ein neuer FB (Freigabe-Lichtschranke-FB24) aufgebaut in SCL. Im Diagramm 3.17 ist der Zusammenhang
zwischen den einzelnen Modi und der Freigabe dargestellt.
FB
Handbetrieb
Signalaustausch
FB24
Freigabe
Lichtschranke
Status / Freigabe
FB
Prozesssteuerung
(Ein-/Auslagerung)
Rückmeldung / Fehler
FB
Störungen
Abbildung 3.17: Einbindung des Bausteins FB24 Freigabe\_Lichtschranke in das
Gesamtsystem
Der Baustein FB24 Freigabe\_Lichtschranke ist ein zentraler Bestandteil der Sicher
heitsarchitektur.
Ohne ihn wäre ein sicheres Abschalten im Fehlerfall nicht möglich. Die Prozesssteue
rung wird durch die Freigabe in jedem Schritt überwacht.
Durch die Einbindung in den FB11 Strungen erfolgt eine Überwachung sämtlicher
Ausgänge über die Variable b\_Freigabe. Diese wird gesetzt, sobald keine Störungen
aktiv sind und die Freigabe durch das Lichtgitter gegeben ist.
Automatisiertes Modell- Hochregallager
31
3 Implementierung eines Lichtgatters in das Gesamtkonzept
Abbildung 3.18: Setzen der Variable b\_Freigabe in FB11 Strungen
Jene Variable in Abbildung 3.18 wurde bereits in vorherigen Projekten verwendet.
Sie wird im Block FB\_Ausgnge eingesetzt, als zusätzliche invertierte Rücksetzbedin
gung. Somit ist ein fehlerfreier Betriebsmodus der einzelnen Aktoren gewährleistet.


\subsubsection{Logische Abschaltbedienungen der Lichtschranke}

In Abbildung 3.19 ist der Ablauf der Sicherheitsabschaltung grafisch dargestellt. Sobald
der Prozesswert NBO größer als 3 ist, wird die Aktorik abgeschaltet. Das bedeutet,
dass eine sofortige Abschaltung erfolgt, sobald zusätzlich zum Zylinder eine Hand oder
ein anderes Objekt detektiert wird– also ein weiterer Lichtstrahl unterbrochen ist.


\subsection{Visualisierung auf dem HMI}
In diesem Kapitel wurde die Implementierung einer sicherheitsgerichteten Lichtschranke
in das bestehende Hochregallager-System beschrieben. Dabei erfolgte zunächst die
Analyse der Prozesswerte des IFM-Lichtgatters sowie die Auswahl des geeigneten
Betriebsmodus zur sicheren Erkennung von Objekten im Gefahrenbereich. Anschlie
ßend wurde der vollständige Systemaufbau erläutert und die Anbindung über den
IO-Link-Master an die Siemens-SPS umgesetzt. Die daraus resultierende Sicherheitslo
gik ermöglicht es, gefährliche Bewegungen im Fehlerfall automatisch zu stoppen und
so die funktionale Sicherheit des Systems zu gewährleisten.
Die Visualisierung der Sicherheitszustände erfolgt über den integrierten Störungsbau
stein des HMI-Panels, wodurch Unterbrechungen und Fehlermeldungen automatisch
angezeigt werden. Eine Erweiterung dieser Funktion um eine gezielte Anzeige des
Anlagenzustands (z.B. „Normalbetrieb“, „Sicherer Halt“) bietet zusätzliches Potenzial
zur Verbesserung der Bedienerfreundlichkeit und Systemtransparenz.
Damit ist die sicherheitstechnische Basis des Hochregallagers geschaffen, auf der im
folgenden Kapitel weitere Funktionen und Erweiterungen aufgebaut werden können

